% ============================================================
% Chapter 4: Testing Methodology
% ============================================================
\chapter{Testing Methodology}
\label{chap:testing}

\section{Derivation of the Verification Strategy}
\label{sec:derivation_of_strategy}

The two-layer verification strategy applied in this project is not a novel invention but a pragmatic synthesis and direct application of established models from functional safety and software engineering. Its structure is specifically tailored to the challenges of a retrofit project, where a legacy function must be verified against modern standards.

The scientific foundation for this approach is the \textbf{V-model}, a cornerstone of safety-critical system development that underpins the safety lifecycle of the foundational functional safety standard \textbf{IEC 61508} and its domain-specific derivatives such as \textbf{EN 50716}. The V-model posits that each phase of specification and design is mirrored by a corresponding verification or validation phase.

Our methodology maps directly to the right side of this V-model:
\begin{itemize}
    \item \textbf{Layer 1 (Component \& Interface Tests)} corresponds to the lower levels of verification in the V-model: the \textbf{unit and integration tests}. This layer proves that the individual design artifacts (VHDL modules, hardware interfaces) are correctly implemented against their technical specifications (timing diagrams, logic requirements). This constitutes a classic white-box verification approach.
    \item \textbf{Layer 2 (System-Level Tests)} corresponds to the highest level of validation in the V-model: the \textbf{system and acceptance tests}. Here, the fully integrated system is validated as a black box against the original, functional system requirements (the behavior of the legacy system).
\end{itemize}

The overarching process structure for planning, executing, and documenting these tests is guided by the V\&V processes defined in the \textbf{IEEE 1012} standard~\cite{ieee1012_2024}. Combining these established models into an application-specific strategy was necessary because no single model adequately covers the dual requirement of verifying new technology (the FPGA) while simultaneously validating it against a legacy system's behavior.

Demonstrating the functional equivalence of the retrofitted anti-skid test-diagnosis module to its legacy counterpart~\cite{rexroth1992wsp}, while ensuring compliance with contemporary safety standards, necessitated a structured and multi-faceted verification and validation (V\&V) strategy. To satisfy the V\&V requirements inherent in standards such as~\gls{iec61508}~\cite{iec61508_2010} and~\gls{ieee1012}~\cite{ieee1012_2024} for this safety-related retrofit, a comprehensive multi-layered testing methodology was implemented. This methodology aimed to provide robust evidence, addressing functional requirements derived from reverse engineering the legacy system~\cite{rexroth1992wsp, ayoub2014, leyer2023} and incorporating principles for verifying modifications to safety systems as outlined in the railway software standard~\gls{en50716}:2023~\cite{en50716_2023}.

Central to this approach is a two-layer verification strategy~–~\textit{Component/Interface} and \textit{System-Level} testing~–~designed to facilitate targeted verification and progressive integration, allowing for effective issue isolation. This core strategy was selected to balance the need for detailed component-level verification (Layer~1) with the practical challenges of system integration testing (Layer~2), providing a pragmatic and traceable path towards demonstrating equivalence. Subsequent work by follow-on teams demonstrated the methodology's value: the Layer~2 fixture was instrumental in uncovering and diagnosing critical hardware-software integration bugs.

% ------------------------------------------------------------
\section{Layer 1: FPGA Logic and Interface Verification (Deterministic)}
\label{sec:fpga_interface_testing}

This foundational layer focused on verifying the internal logic of the~\gls{fpga} (\Cref{sec:fpga_core_details}) and its immediate electrical interfaces using deterministic, white-box methods. The primary goals were to confirm logical correctness, validate timing accuracy against expected behaviour derived from~\gls{mcs48} documentation and legacy system analysis~\cite{rexroth1992wsp}, and ensure basic hardware integrity before system integration. This approach aligns with the hardware verification requirements of~\gls{en50129} (covering the HDL/FPGA implementation) and the software requirements of~\gls{en50716}, and draws on established structured testing principles~\cite{watson1996}.

\begin{description}
	\item[VHDL Simulation and Testbenches:] Extensive pre-synthesis and post-synthesis simulations using the simulator bundled with the Microsemi Libero SoC toolchain (\textbf{ModelSim ME}, the Mentor/Siemens ModelSim Microsemi Edition) verified the~\gls{fpga} design prior to hardware deployment. Dedicated~\gls{vhdl} testbenches (e.g., \texttt{atop\_tb.vhd}, see~\Cref{sec:app_vhdl_testbenches}), adapted from prior work~\cite{leyer2023} and expanded, provided controlled and repeatable stimuli to the Device Under Test (DUT)~– the top-level~\gls{fpga} design and its sub-modules. Key functionalities verified included: the logical behaviour of the emulated~\gls{mcs48} core by running test instruction sequences; memory interface simulations confirmed that timing parameters (setup, hold, access times) for accessing the external Flash~\gls{rom} (\texttt{S29AL016J}) and~\gls{nvram} (\texttt{CY14B104NA}) met the requirements specified in the respective component datasheets; correct operation of emulated peripheral logic (7-segment display latching, button input debouncing/interrupt logic, frequency sweep generation state machine, \gls{spi}~master communication protocol); and the crucial watchdog interaction logic (correct \texttt{!INT} pulsing, proper response to~\texttt{!FPGA\_8048\_RESET}). Assertions within the testbenches and manual waveform analysis (\Cref{sec:app_verification_captures}, e.g.,~\Cref{fig:watchdog_reaction}) confirmed expected behaviour under various stimuli. These simulation results were critical for informing iterative corrections to the~\gls{hdl} code (\Cref{subsec:fpga_interface_results}) and establishing baseline confidence in the logic design.
	\item[Boundary Scan (JTAG) Checks:] \gls{jtag} (\gls{ieee1149}) was used for basic infrastructure verification: confirming~\gls{fpga} ID (\texttt{A3P1000}), successful programming of the bitstream, and performing initial I/O pin connectivity checks to detect manufacturing defects like shorts or opens between adjacent pins accessible via the boundary scan chain.
	\item[Hardware Timing Measurements:] An oscilloscope was used to validate critical timing parameters on the assembled Rev~1.0 hardware, specifically aspects sensitive to physical implementation: watchdog pulse periodicity and frequency sweep accuracy (verifying the generated ramp profile from \SI{850}{\hertz} to \SI{1550}{\hertz} and its duration). These results are presented in~\Cref{subsec:fpga_interface_results}.
\end{description}
This deterministic layer established confidence in the basic functional correctness and timing behaviour of the isolated~\gls{fpga} logic and its direct interfaces, providing a necessary foundation for the integrated system testing described next.

% ------------------------------------------------------------
\section{Layer 2: Integrated System Scenario Testing (Fixture-Based)}
\label{sec:system_scenario_testing}

\begin{figure}[htbp]
	\centering
	\includegraphics[width=0.7\textwidth]{ayoub_test_setup}
	\caption[Conceptual Test Fixture Setup]{Conceptual overview of the Layer 2 test fixture setup, illustrating the principle of stimulating and monitoring the Device Under Test (DUT). Based on the original test concept from Ayoub~\cite{ayoub2014}.}
	\label{fig:layer2_setup}
\end{figure}

Building upon component verification, this layer validated the integrated system~– comprising the Base Board with the programmed~\gls{fpga}, the Upper Board, the~\gls{stm32} Diagnostics Board, and the LED Display Board~– under operational scenarios using a dedicated hardware test fixture (\Cref{fig:layer2_setup}). This aimed to demonstrate functional equivalence with the legacy system's behaviour in key operational modes and fault conditions, adopting a black-box perspective at the board level. This fixture-based setup also enabled a follow-on team to diagnose several deep-seated hardware-software integration bugs.

\begin{description}
	\item[Dedicated Test Fixture (`Testkartentester`):] Testing relied significantly on the hardware test fixture developed by Dolata et al.~\cite{dolata2023}. This fixture provides the necessary physical connection points (via the main system connector, \texttt{X4}) to the retrofitted board assembly and interfaces with a PC via an intermediate~\gls{stm32} Nucleo board (\texttt{STM32F401RE}, see~\cite[Abschn.~6.2]{dolata2023}). It allows the PC-based test software to stimulate the board's inputs (e.g., simulating button presses or specific signal conditions corresponding to legacy faults using pins like \texttt{18z}, \texttt{24d}) and monitor its outputs (e.g., reading display states indirectly or checking signal lines).
	\item[YAML-Based Scenario Definition:] To ensure structured and repeatable system-level testing, test procedures were defined in a human-readable~\gls{yaml} format (\Cref{sec:app_yaml_test_cases}, see files \texttt{tests.yml}, \texttt{tests\_all.yml} in context). Each test case specified the sequence of input stimuli (\texttt{pin\_sets}) to be applied via the fixture, the expected system state (primarily the codes shown on the 7-segment displays, \texttt{expected\_err\_code}), expected stored fault sequences (\texttt{expected\_err\_seq}, read via the STW button simulation), and prompts for manual tester interaction or verification where necessary (e.g., confirming visual display, pressing the simulated \enquote{Lösch} button via the \texttt{del: True} flag). This approach addresses requirements coverage by explicitly mapping test scenarios to legacy system functions and fault codes documented in~\cite{rexroth1992wsp}. For example, simulating a \enquote{Geber 3} fault involves setting \verb|pin_sets: {18z: 1, 24d: 0}| and expecting error code \texttt{"31"} and sequence \texttt{"03 31 09"}.
	\item[Interactive Test Execution:] A Python application (\texttt{main.py}, leveraging modules like \texttt{run\_tests.py} and \texttt{load\_tests.py}, see \texttt{Testkartentester} context) running on a PC orchestrated the test execution based on the loaded~\gls{yaml} scenarios. It communicated with the test fixture (via the Nucleo board using \texttt{pyserial}) to apply the defined stimuli, paused execution using a text-based UI (\texttt{gui/testing.py}) to prompt the tester for required manual actions or observations (e.g., \enquote{Press STW Button}, \enquote{Verify Display shows 31}), and logged the pass/fail outcome based on the tester's input (\texttt{logger/test\_logger.py}, see~\Cref{sec:app_system_log_excerpts}). This semi-automated approach allowed for efficient execution of complex scenarios while incorporating necessary human checks for outputs like visual displays.
	\item[Key Functionality Validated:] This fixture-based approach enabled end-to-end validation of critical legacy functionalities, directly addressing functional equivalence requirements outlined in~\cite{rexroth1992wsp}:
	      \begin{itemize}
		      \item \textit{Complete Fault Memory Cycle:} Triggering specific fault conditions via input pins (e.g., Geber~3 fault via pins \texttt{18z}/\texttt{24d}), verifying the correct immediate error code display (\texttt{31}), simulating the \enquote{STW} button press to read the stored error sequence (\texttt{03 31 09}), verifying correct storage and recall across simulated power cycles (managed by the test script controlling fixture power), and confirming successful clearing of faults using the simulated \enquote{Lösch} button (resulting in the \enquote{99} display).
		      \item \textit{Self-Test Activation and Execution:} Simulating the \enquote{Test} button press via the fixture to initiate the legacy self-test sequence, verifying the correct sequence of codes displayed on the 7-segment displays (\enquote{88}, \enquote{85}, etc.), and confirming the activation of the frequency sweep output during the appropriate phase (timing measured in Layer~1).
		      \item \textit{User Interface Equivalence:} Verifying that the behaviour of the system in response to simulated button presses (Test, STW, Lösch, Tür/V) and the resulting display states correspond directly to the behaviour described in the original system manual~\cite{rexroth1992wsp}.
		      \item \textit{Integrated System Stability:} Running sequences of tests and power cycles (e.g., from \texttt{tests\_all.yml}) to check for unexpected hangs or failures in the integrated system.
		      \item \textit{Diagnostics Subsystem Interaction:} Verifying that diagnostic messages were correctly transmitted from the~\gls{fpga} to the~\gls{stm32} board via the isolated~\gls{spi} link (\Cref{sec:stm32_details}) during fault scenarios and logged appropriately on the~\gls{sd}~card (verified post-test, see~\Cref{subsec:system_scenario_results}).
	      \end{itemize}
\end{description}
These fixture-based tests were crucial for demonstrating functional equivalence under relevant operational conditions~\cite{losada2021}, providing system-level evidence complementing the component-level verification of Layer~1.

% ------------------------------------------------------------
\section{Validation Context and Complementary Approaches}
\label{sec:validation_context}

The primary two-layer testing strategy was situated within a broader validation context, acknowledging both the project's scope and the requirements for safety-critical systems.

\begin{description}
	\item[\gls{hil} Principles Alignment:] While a full dynamic vehicle model interacting with the hardware in real-time was beyond the project's scope and resources, the developed methodology aligns with key~\gls{hil} principles~\cite{hil_rail_suspension_2023, zhou_wu_2020}. Specifically, it combines detailed component-level simulation (Layer~1) with closed-loop functional testing in a controlled hardware environment (the Layer~2 fixture testing). This allows for the execution of realistic operational scenarios and fault conditions, covering aspects typically addressed in industry~\gls{hil} campaigns~\cite{wabtec_wsp_cert_2022}, and provides a foundation for future integration into a more comprehensive~\gls{hil} test rig (\Cref{sec:future_work}).
	\item[Fault Handling Verification:] The focus of fault testing in Layer~2 was on verifying the system's correct detection and response to known legacy error conditions (as defined in~\cite[Blatt~20/25]{rexroth1992wsp} and triggered via fixture inputs defined in files like \texttt{tests.yml}). This aligns directly with the primary goal of demonstrating functional equivalence. Exhaustive low-level hardware fault injection (e.g., simulating stuck-at faults or bit-flips within the~\gls{fpga} using techniques surveyed in~\cite{ziade2004faultinjection}) was not performed, as the main objective was replicating existing behaviour rather than proving resilience against arbitrary internal hardware faults beyond the scope of the original design's fault detection mechanisms.
	\item[Formal Verification Scope:] Verification relied primarily on simulation (Layer~1) and testing (Layer~1 and Layer~2). Formal methods, such as model checking or formal equivalence proving~\cite{formal_verification_rail_2023, kravets2005formality}, were not applied to the~\gls{hdl} or software components due to project scope and complexity. However, rigorous code reviews and structured~\gls{vhdl} testbenches aimed to ensure logical correctness and address potential design flaws early in the process.
	\item[Test Coverage Considerations:] The verification strategy aimed for comprehensive \textit{functional coverage} based on the legacy system's requirements and documented behaviour~\cite{rexroth1992wsp}. Layer~1 simulations targeted structural coverage of the~\gls{hdl} code through directed testbenches, aiming to exercise critical logic paths. Layer~2 scenario testing aimed for requirements-based coverage, ensuring all documented fault codes, self-test sequences, and user interactions were validated. While quantitative code coverage metrics (e.g., statement/branch for C, structural coverage for VHDL) were not systematically measured due to tool limitations or scope, the methodology focused on achieving high confidence through targeted functional and scenario-based validation. Structured testing principles~\cite{watson1996} influenced the design of test cases to maximize the testing of independent logic paths within the system's core functions.
	\item[Environmental Qualification Approach:] Compliance with environmental standards for rolling stock equipment (\gls{en50155}:2021~\cite{en50155_2021} for operational conditions and~\gls{emc} via~\gls{en50121}~\cite{en50121_3_2_2016}; \gls{en61373}:2010~\cite{en61373_2010} for shock and vibration) currently relies on the design choices made following established principles (\Cref{sec:design_principles}). Formal demonstration of compliance would require accredited laboratory testing on the finalised hardware, which is considered a necessary step for full certification but was outside the scope of this research phase.
\end{description}
These complementary approaches help define the boundaries and context of the verification performed, highlighting the specific focus on functional equivalence demonstrated through simulation and fixture-based testing, while acknowledging areas requiring further work (like accredited environmental testing) for full certification.

% ------------------------------------------------------------
\section{Verification Artefacts and Process Control}
\label{sec:verification_artefacts}

Generating auditable evidence suitable for supporting a safety case according to standards like~\gls{iec61508}~\cite{iec61508_2010}, \gls{en50129}~\cite{en50129_2018}, \gls{en50716}~\cite{en50716_2023}, and~\gls{ieee1012}~\cite{ieee1012_2024} was a key process objective throughout the verification activities.

\begin{description}
	\item[Test Documentation and Logs:] Core artefacts generated and managed include:
	      \begin{itemize}
		      \item \gls{vhdl} Testbenches (\Cref{app:test_cases}): Providing documented, executable specifications for Layer~1 verification (e.g., \texttt{atop\_tb.vhd}).
		      \item Simulation Waveform Captures (\Cref{app:test_logs}): Visual evidence of signal behaviour during Layer~1 simulations (e.g.,~\Cref{fig:watchdog_reaction}).
		      \item \gls{yaml} Test Definitions (\Cref{app:test_cases}): Documented scenarios for Layer~2 system testing (e.g., \texttt{tests.yml}, \texttt{tests\_all.yml}).
		      \item Fixture Test Logs (\Cref{app:test_logs}): Detailed records of executed Layer~2 tests, including tester inputs and pass/fail results (e.g., \texttt{test\_results.log}, see~\Cref{sec:app_system_log_excerpts}).
		      \item Hardware Measurement Records: Oscilloscope screenshots and documented measurements for Layer~1 timing verification.
	      \end{itemize}
	\item[Traceability:] While formal, fine-grained traceability matrices were not generated within this project's scope, traceability, a key requirement of~\gls{iec61508} and~\gls{ieee1012}, was maintained implicitly through the structured testing approach. Layer~1 tests trace to~\gls{hdl} module requirements (derived from RE of the~\gls{mcs48} peripherals and memory interfaces), while Layer~2 YAML scenarios trace explicitly to legacy functional requirements documented in~\cite{rexroth1992wsp} (e.g., specific fault codes, self-test sequences, button functions). The verification results (\Cref{chap:results}) link these tests back to the design elements (\Cref{chap:architecture}). For formal certification, these implicit links would need to be captured in explicit traceability matrices demonstrating complete coverage from requirements to verification results. The checklist in~\Cref{chap:safety_checklist} provides a high-level mapping as a step towards this.
	\item[Version Control:] Git was rigorously used for configuration management of all design files (\gls{hdl}, schematics, layouts), test artefacts (\gls{vhdl}~testbenches, \gls{yaml}~scenarios), firmware (\gls{stm32}~C code), and documentation (LaTeX source, reports), ensuring reproducibility and change tracking as required by~\gls{en50716}. This provides an audit trail for all modifications and verification activities. The \texttt{.gitignore} file ensures build artifacts are correctly excluded.
	\item[Tool Qualification Considerations:] Standard, widely used industry tools (Microsemi Libero SoC including ModelSim ME, STM32CubeIDE, Python with libraries like \texttt{pyserial}) were employed. Formal certification according to~\gls{iec61508} or~\gls{en50716} would necessitate a tool qualification assessment. This involves classifying each tool (e.g., T1, T2, T3 based on potential impact) and providing justification for its use, potentially including vendor qualification data, validation suites, or evidence of operational experience, depending on the required Tool Confidence Level (TCL). This assessment was beyond the current project scope but is acknowledged as a required certification activity.
\end{description}

\section{Compliance with EN 50716 for the Verification of Modifications}
\label{sec:en50716_modifications}

A retrofit is, by definition, a modification to an existing system. The EN 50716 standard specifies requirements for managing such changes to ensure that system safety is not compromised. Our verification methodology was deliberately designed to fulfill these core requirements.

\begin{description}
    \item[Impact Analysis:] An implicit impact analysis was performed prior to implementation. It was clearly defined that the modification was confined to the hardware of the test-diagnosis board, while the safety-critical firmware of the emulated controller remained untouched. This precisely defined the scope of the necessary verification.

    \item[Verification of the Change:] The entire two-layer methodology serves as the formal process for verifying the modification. Layer 1 verifies the correct implementation of the new hardware and HDL components at a technical level. Layer 2 verifies that the change correctly implements the required system behavior.

    \item[Regression Testing:] The Layer 2 scenario tests constitute a comprehensive regression test suite. By systematically re-validating all known legacy functions (fault memory, self-tests), it is ensured that the modification has not introduced any unintended adverse effects on existing, functional parts of the system.

    \item[Configuration Management:] As detailed in \Cref{sec:verification_artefacts}, the consistent use of Git for all artifacts (hardware design, HDL, software, test cases, documentation) ensured complete traceability and reproducibility of all changes, which is a central requirement of the standard.
\end{description}

This systematic approach, combining targeted verification layers with defined process controls, provides the auditable evidence base discussed and analysed in the following chapter.
